For $F=\C((t))$, the conclusion over the ground base may be
transported by the already constructed comparison $\Theta$:
it becomes $\Vol([A])=\Vol([A'])$.
The hypothesis still concerns every $b\in B$ in the
saturated ambient model, with its own coefficient base
$F\langle b\rangle$.
It is not enough to test only the points of $B$ in the
Puiseux field.

For arbitrary $b$, retain the intrinsic ring
$R(F\langle b\rangle)$.
An algebraically closed ambient valued field does not enlarge
the definability base.
In particular, adjoining an algebraic closure of the residue
field to the coefficient field of a realization can lose
arithmetic descent; an action recording only tame inertia
does not restore it.

Here is a concrete instance.
Choose a unit $b$ with residue $s$ transcendental over $\C$
and put $K=F(b)$ with its Gauss valuation.
For a polynomial $\sum_i a_i b^i$, its terms of least
$t$-valuation have nonzero residual polynomial in $s$.
Consequently
\[
 k(K)=k=\C(s),\qquad \Gamma(K)=\Z.
\]
The residue set
\[
 L=\{z:z^2=s\}
\]
is defined over $K$.
It is a nonsplit quadratic residue cover: $s$ is not a square
in $\C(s)$, since its order at $s=0$ is odd.
Over $\overline{k}$ it becomes two points, each with trivial
tame-inertia action, just as the constant two-point set does.

The distinction survives in the RES quotient.
Indeed, at this Gauss base the same finite framing argument
used for $\Theta$ normalizes every finite rational value
support by roots of $t$.
RV quantifier elimination gives constructible sets and graphs
over $k$, with a finite cyclic inertia action and with their
$k$-descent retained.
Changing a framing root multiplies coordinates by roots of
unity in $\C$, so it gives a $k$-defined isomorphism after
forgetting inertia.
Applying compactly supported \'etale Euler realization then
gives an additive and multiplicative invariant
\[
 \rho:R(K)\longrightarrow
 K_0\bigl(\operatorname{Rep}_{\Q_\ell}(G_k)\bigr),
 \qquad G_k=\operatorname{Gal}(\overline{k}/k).
\]
Here the realization is
$\sum_j(-1)^j[H_c^j(-_{\overline{k}},\Q_\ell)]$;
its additivity and multiplicativity follow from the
localization sequence and K\"unneth formula.\footnote{For
the standard realization on $K_0(\var_k)$, see
J.~Nicaise, \emph{A trace formula for varieties over a
discretely valued field}, Section~2.1(2),
\href{https://arxiv.org/pdf/0805.1323}{arXiv:0805.1323}.
Only this standard cohomological realization is used in this
diagnostic; no comparison theorem for RES is assumed.}
The quotient relations are killed because both
$q^\sharp$ and $\G_m$ normalize to $\G_m$ after inertia is
forgotten. Thus this argument constructs the required
invariant without asserting a new comparison isomorphism.
If $\varepsilon$ is the quadratic character attached to
$k(\sqrt{s})/k$, then
\[
 \rho([L])=[1]+[\varepsilon]\ne 2[1].
\]
For $g\in G_k$ with $\varepsilon(g)=-1$, the two traces
are $0$ and $2$. Residue base change gives two inertia-fixed
points in both cases, so the proposed comparison is not
injective.

There is also no restriction of the theorem's test to points
generating unramified extensions of $F$.
For example, take
\[
 B=\{b:b^2=t\},\qquad A=B,\qquad f=\mathrm{id}_B,
 \qquad A'=\varnothing.
\]
Every $b\in B$ has value $1/2$, so $F(b)/F$ is ramified
quadratic. The unramified-point test has no points to check.
On the other hand, $\rv$ identifies $B$ with the finite
RES object
\[
 T=\{u:u^2=\rv(t)\}\in\RES[0].
\]
The two roots have distinct RV classes; conversely these are
all the roots of $u^2=\rv(t)$.
The degree-zero lifting comparison therefore gives
\[
 \Vol([B])=\Theta([T])
   =[\{x:x^2=1\},\ \zeta\cdot x=\zeta x].
\]
The action factors through $\mu_2$ and interchanges the two
points. Forgetting the action and taking compactly supported
Euler characteristic sends this class to $2$, so it is
nonzero. This Euler characteristic respects the linear-action
relations, since every affine space of a fixed dimension has
Euler characteristic $1$.
Thus the asserted unramified-only criterion would give a
false conclusion. The theorem retains the test at all
ambient parameters in their original RES coefficient rings.
